% Lösungen Einheit 5 von 5 – Körper im Koordinatensystem und Drehungen (zusammenbau v0.1)
\einheitenkopf{Einheit 5 von 5 · Körper im Koordinatensystem und Drehungen}

\erg{19}{a) Sie liegt parallel zur $x_1x_2$-Ebene in der Höhe $4$. \quad b) $F$ liegt senkrecht über $C$ in der Höhe von $D$: $F(2 | 4 | 3)$ (nicht mit $E$ verwechseln). \quad c) Z. B. rechter Winkel im Ursprung, die Katheten auf die Achsen: $(0 | 0 | 0)$, $(6 | 0 | 0)$, $(0 | 8 | 0)$, Spitze $(0 | 0 | 5)$ – die Hypotenuse $10$ gehört nicht auf eine Achse. \quad d) Nur die Ecke $(0 | 2 | 2)$ passt: danach $(-1 | 1 | 3)$. \quad e) Alle Grundecken haben dieselbe dritte Koordinate: Die Grundfläche ist parallel zur $x_1x_2$-Ebene. Höhe $h = 7 - 2 = 5$ (nicht $7$). \quad f) $P$ liegt in der $x_1x_2$-Ebene. Mittelpunkt der Grundfläche $N(4 | 3 | 0)$ (nicht $M$): $|\overrightarrow{NP}| = \sqrt{9 + 16} = 5$ – $P$ liegt auf dem Rand. \quad g) Strahlensatz am Querschnitt, von der Traufe aus: $d = \frac{2 - 0{,}5}{5{,}5 - 0{,}5} \cdot 4 = 1{,}2$ m (Höhendifferenz, nicht $\frac{2}{5{,}5}$). Nicht zählend: zwei Streifen, Anteil $\frac{2 \cdot 1{,}2}{8} = 30\,\%$. \quad h) $B$ läuft auf einem Kreis um die $x_3$-Achse mit Radius $10$; die dritte Koordinate bleibt $1$. Z. B. $(6 | 8 | 1)$, $(6 | -8 | 1)$, $(8 | 6 | 1)$. \quad i) $C$ läuft auf einem Kreis um den Lotfußpunkt $L(4 | 4 | 0)$ auf $AB$ (nicht um den Ursprung), Radius $|\overrightarrow{LC}| = \sqrt{18}$. In der $x_1x_2$-Ebene senkrecht zu $AB$: $L \pm 3 \cdot (1 | 1 | 0)$; mit $x_1 < 4$: $C'(1 | 1 | 0)$.}
\erg{20}{a) Fehler: $(6 | 2 | 4)$ liegt über $B$, das ist $E$. Richtig: $F$ liegt über $C$, also $F(1 | 5 | 4)$. \quad b) Gerade heißt: Die Seitenkanten stehen senkrecht auf der Grundfläche, also parallel zur $x_3$-Achse. Entlang dieser Richtung ändert sich nur $x_3$.}
